\chapter[Introdução]{Introdução}

Um drone em patrulhamento visita pontos de interesse de forma recorrente com deslocamento mínimo. Em inspeção, vigilância ambiental e monitoramento de infraestrutura, a qualidade de uma rota não se reduz à distância total percorrida. Curvas bruscas aumentam o tempo de voo, o consumo de energia e a dificuldade de execução da trajetória. Tratar cada deslocamento como aresta independente entre dois pontos desconsidera esse custo.

O \ac{TSP} modela rotas que partem de uma origem, visitam todos os pontos exatamente uma vez e retornam à origem \cite{applegate2006traveling}. Como a versão de decisão do \ac{TSP} é \ac{NP}-completa \cite{garey1979computers}, enumerar todas as rotas torna-se inviável conforme o número de pontos cresce. Usamos heurísticas e meta-heurísticas para obter soluções de baixo \textit{makespan} em tempo computacional controlado.

Estudamos uma variante do \ac{TSP} motivada por patrulhamento aéreo com drones: o \ac{TSP-SD-ATP}. Nessa formulação, o custo de seguir de um ponto atual para um próximo ponto depende também do ponto visitado imediatamente antes. A instância usa um tensor tridimensional de custos, indexado por nó anterior, nó atual e próximo nó, em vez de uma matriz bidimensional de arestas. O objetivo é minimizar o \textit{makespan}, ou seja, o tempo/custo total da rota da base até o retorno.

\section{Motivação}

Pesquisadores têm estudado o roteamento com drones em diferentes formulações. Murray e Chu \cite{murray2015flying} introduziram o \textit{Flying Sidekick Traveling Salesman Problem}, no qual caminhão e drone cooperam em entregas. Agatz, Bouman e Schmidt \cite{agatz2018optimization} estudaram o \ac{TSP} com drone por meio de formulações e heurísticas, enquanto Dell'Amico et al. \cite{dellamico2021multiple,dellamico2022exact} avançaram variantes com múltiplos drones e métodos exatos. No contexto de patrulha, Rajan, Sundar e Gautam \cite{rajan2022routing} modelaram missões de \acp{VANT} com horizonte finito e decisões recorrentes de visita.

Pequenas alterações operacionais geram variantes distintas do \ac{TSP}. Este trabalho não combina caminhão e drone nem modela múltiplos veículos; isolamos o caso de um único drone que deve visitar todos os pontos de interesse e retornar à base. A contribuição está em avaliar o impacto de uma função de custo dependente da sequência, na qual a transição anterior--atual--próximo incorpora distância euclidiana e penalidade angular.

Haroun et al. \cite{haroun2015performance} e Chandra e Naro \cite{chandra2022comparative} comparam \ac{GA}, \ac{PSO} e \ac{ACO} no \ac{TSP} clássico e em variantes próximas. Nenhum desses trabalhos compara sistematicamente os três métodos no \ac{TSP-SD-ATP}.

\section{Objetivos e Desafios da Pesquisa}

O objetivo geral deste trabalho é analisar o desempenho de meta-heurísticas bioinspiradas na otimização de rotas de patrulhamento com drones modeladas como \ac{TSP-SD-ATP}.

Os objetivos específicos são:
\begin{alineas}
  \item formalizar a função objetivo baseada em distância euclidiana e penalidade angular dependente da sequência;
  \item implementar e comparar \ac{GA}, \ac{PSO} e \ac{ACO} sobre a mesma representação de instância e a mesma função objetivo;
  \item usar busca exaustiva como referência ótima nas instâncias pequenas em que a enumeração completa foi executada;
  \item implementar um limitante inferior por relaxação do problema de designação para apoiar a análise em instâncias maiores;
  \item medir qualidade da solução, tempo de otimização, estabilidade entre sementes e comportamento de convergência;
  \item discutir limitações metodológicas associadas a instâncias sintéticas, parâmetros fixos e diferenças de codificação entre métodos.
\end{alineas}

A complexidade combina duas fontes. A primeira é combinatória: a quantidade de rotas cresce fatorialmente com o número de pontos. A segunda é estrutural: o custo de uma transição depende de três nós consecutivos, não apenas da aresta atual. Métodos projetados para o \ac{TSP} clássico exigem adaptação para avaliar rotas em um tensor de custos tridimensional.

\section{Hipótese}

A hipótese principal é que meta-heurísticas bioinspiradas produzem rotas factíveis de baixo \textit{makespan} em instâncias nas quais a busca exaustiva se torna inviável. A hipótese secundária é que a qualidade e o tempo de execução variam substancialmente conforme a forma como cada método representa e explora a dependência de sequência do \ac{TSP-SD-ATP}.

As perguntas de pesquisa são:
\begin{alineas}
  \item Nas instâncias pequenas, qual é o gap das meta-heurísticas em relação ao ótimo obtido por busca exaustiva?
  \item Nas \TotalInstancias{} instâncias avaliadas, qual método apresenta menor \textit{makespan} mediano?
  \item Qual é o custo temporal de obter soluções de melhor qualidade?
  \item As diferenças observadas entre \ac{GA}, \ac{PSO} e \ac{ACO} são estatisticamente consistentes quando os métodos são comparados por instância?
  \item O limitante inferior por relaxação do problema de designação é informativo para esta variante com penalidade angular?
\end{alineas}

\section{Contribuições}

Este trabalho contribui com:
\begin{alineas}
  \item a formulação computacional do \ac{TSP-SD-ATP} para rotas de patrulhamento com drones, usando tensor 3D de custos;
  \item uma implementação em Go de \ac{GA}, \ac{PSO}, \ac{ACO}, busca exaustiva e \ac{LB} por relaxação \ac{AP}/Hungarian sobre a mesma função objetivo;
  \item uma base experimental com \TotalExecucoes{} arquivos de resumo, evolução e tempo, cobrindo \TotalInstancias{} instâncias e \TotalSementesGA{} sementes por meta-heurística;
  \item a comparação com ótimo exato em \TotalInstanciasComOtimo{} instâncias pequenas e com \ac{LB} em todas as \TotalInstancias{} instâncias;
  \item uma análise empírica e estatística do trade-off entre qualidade de solução e tempo computacional.
\end{alineas}

\section{Organização da Monografia}

O \autoref{fundamentacao} apresenta a fundamentação teórica sobre \ac{TSP}, custos dependentes da sequência, roteamento com drones e meta-heurísticas bioinspiradas. O \autoref{proposta} descreve a formulação adotada, a arquitetura da implementação, os métodos comparados, o \ac{LB} e o pipeline experimental. O \autoref{experimentos} apresenta a configuração dos experimentos, os resultados, os testes estatísticos e a discussão das limitações. O \autoref{conclusao} sintetiza as conclusões, contribuições e trabalhos futuros.
